\documentclass[11pt,a4paper]{article}

% ---- Packages ----
\usepackage[utf8]{inputenc}
\usepackage[T1]{fontenc}
\usepackage{geometry}
\geometry{margin=2.5cm}
\usepackage{amsmath,amssymb}
\usepackage{graphicx}
\usepackage{booktabs}
\usepackage{caption}
\usepackage{subcaption}
\usepackage{hyperref}
\usepackage{natbib}
\usepackage{setspace}
\usepackage[table]{xcolor}
\usepackage{float}
\usepackage{csquotes}
\usepackage{siunitx}
\usepackage{lscape}

\onehalfspacing

\sisetup{group-separator={,}, group-minimum-digits=5}

\hypersetup{
    colorlinks=true,
    linkcolor=blue,
    citecolor=blue,
    urlcolor=blue
}

\newcommand{\Rtwo}{\ensuremath{R^2}}

\begin{document}

%% ======================== TITLE ========================
\begin{titlepage}
\centering
\vspace*{3cm}

{\LARGE \textbf{Estimating Poverty in Niger with Landsat 7 Satellite Imagery, Deep Learning, and Fay--Herriot Small Area Estimation}}\\[1.5cm]

{\large A Reproducible Open-Source Pipeline for Subnational Poverty Mapping}\\[2cm]

\vfill

{\large \today}\\[1cm]

\end{titlepage}

%% ======================== ABSTRACT ========================
\begin{abstract}
Detailed subnational poverty estimates are critical for targeting development interventions, but household surveys in low-income countries often lack the sample size to produce reliable estimates at the administrative level where policy decisions are made. We propose a two-stage pipeline that combines freely available Landsat 7 satellite imagery, a ResNet-18 deep neural network, and Fay--Herriot small area estimation (SAE) to predict mean wealth across 64 departments in Niger using 2012 DHS data. Unlike existing studies that rely on random train--test splits, we employ group 5-fold spatial cross-validation by department to obtain an honest estimate of generalization performance. Our CNN achieves an out-of-fold \Rtwo{} of 0.482 (CV mean \Rtwo{} = 0.320 $\pm$ 0.156). When CNN predictions are incorporated into a Fay--Herriot model alongside urbanization data, the resulting coefficient of variation (CV) is 11.1\% (95\% bootstrap CI: [7.4\%, 15.2\%]), compared to 17.8\% for a model using urbanization alone. These results demonstrate that combining deep learning with formal small area estimation substantially improves poverty predictions and provides reliable uncertainty quantification. All code is openly available as Kaggle and Colab notebooks, enabling reproduction in any country with DHS and Landsat data.
\end{abstract}

\newpage

%% ======================== 1. INTRODUCTION ========================
\section{Introduction}
\label{sec:intro}

Accurate poverty estimates at the subnational level are essential for the effective allocation of development resources, the design of targeted interventions, and the monitoring of progress toward the Sustainable Development Goals. In many low-income countries, however, nationally representative household surveys such as the Demographic and Health Surveys (DHS) provide reliable estimates only at the national or regional level, with insufficient sample size for smaller administrative units such as departments or districts \citep{elbers2003}.

Niger is a compelling case study. Ranked near the bottom of the Human Development Index, the country faces significant data challenges: its 2012 DHS covers approximately 476 clusters across 64 departments, with some departments having as few as one cluster. Traditional direct survey estimates for such small domains have unacceptably large standard errors, rendering them of limited use for policy.

Recent advances in deep learning and satellite remote sensing have opened new avenues for poverty mapping. \citet{jean2016} demonstrated that convolutional neural networks (CNNs) trained on high-resolution satellite imagery can predict household wealth across several African countries, achieving \Rtwo{} values of 0.55 to 0.75. Subsequent work by \citet{yeh2020}, \citet{deepwealth2024}, and others has extended this approach using Landsat imagery, nightlights, and point-of-interest data. However, these studies share a common methodological limitation: they evaluate model performance using random train--test splits, which fail to account for the spatial structure of the data. When clusters from the same administrative unit appear in both training and test sets, the estimated \Rtwo{} may be optimistically biased.

Furthermore, most satellite-based poverty mapping studies stop at the CNN prediction stage, without incorporating predictions into a formal small area estimation framework that quantifies uncertainty at the domain level. This limits their practical utility for policy makers, who require not only point estimates but also measures of reliability.

In this paper, we address both limitations. We propose a two-stage pipeline that:
\begin{enumerate}
    \item Trains a ResNet-18 on 476 Landsat 7 patches with group 5-fold spatial cross-validation by department, producing out-of-fold CNN predictions;
    \item Feeds these predictions into a Fay--Herriot area-level model \citep{fay1979} along with auxiliary variables (urbanization percentage, NDVI), yielding small area estimates with associated standard errors and coefficients of variation.
\end{enumerate}

Our contributions are:
\begin{itemize}
    \item A rigorous spatial cross-validation framework that provides honest generalization estimates, revealing that random splits overstate \Rtwo{} by approximately 0.2.
    \item An ablation study comparing five Fay--Herriot specifications, showing that CNN predictions reduce the coefficient of variation from 17.8\% to 11.1\%.
    \item Bootstrap confidence intervals (1000 iterations) for the FH coefficient of variation, providing uncertainty quantification for model comparisons.
    \item A fully open-source, reproducible pipeline deployable on Kaggle and Google Colab with no proprietary data or software required.
\end{itemize}

\noindent The remainder of this paper is organized as follows. Section~\ref{sec:lit} reviews related work. Section~\ref{sec:data} describes the data. Section~\ref{sec:method} details the methodology. Section~\ref{sec:results} presents results. Section~\ref{sec:discussion} discusses implications and limitations, and Section~\ref{sec:conclusion} concludes.

%% ======================== 2. LITERATURE REVIEW ========================
\section{Related Work}
\label{sec:lit}

The intersection of satellite imagery, machine learning, and poverty measurement has grown rapidly since the landmark study of \citet{jean2016}, who used a CNN pre-trained on ImageNet to extract features from high-resolution satellite and Google Street View images across five African countries. Their transfer learning approach achieved \Rtwo{} values of 0.55 to 0.75 when predicting mean household wealth at the village level, depending on the specification.

\citet{perez2017} specifically used Landsat 7 imagery with a ResNet-18 architecture, achieving an \Rtwo{} of approximately 0.58 across six African countries. \citet{yeh2020} extended this to approximately 20,000 African villages using a dual-input model combining daytime multispectral imagery and nighttime lights, reporting \Rtwo{} values in the 0.60--0.72 range.

Most recently, \citet{deepwealth2024} proposed an open-source deep learning framework trained on 24 African countries and validated on Madagascar, Brazil, and Japan, with a reported \Rtwo{} of 0.69. \citet{putri2023} achieved an \Rtwo{} of 0.84 at the district level in East Java, Indonesia, by combining multi-source satellite imagery with point-of-interest data.

Despite these impressive results, a critical methodological gap persists across all these studies: they rely on random train--test splits that do not respect spatial structure. When clusters from the same administrative area are present in both training and test sets, the model effectively evaluates its ability to interpolate within known areas rather than generalize to unseen ones. \citet{meyer2018} and \citet{karasiak2022} have shown that random cross-validation can dramatically overstate model performance in geospatial applications.

On the small area estimation side, Fay--Herriot models \citep{fay1979} have been widely applied for poverty mapping in official statistics \citep{rao2015, you2021}. However, the integration of CNN-predicted wealth scores as predictors in Fay--Herriot models remains rare. Notable exceptions include \citet{marhuenda2017} and \citet{ratledge2021}, but neither provides a fully open-source pipeline with spatial cross-validation and bootstrap uncertainty quantification.

Our work bridges these two literatures by combining state-of-the-art computer vision methods with formal SAE, while addressing the methodological limitations identified above.

%% ======================== 3. DATA ========================
\section{Data}
\label{sec:data}

\subsection{Survey Data}
We use the 2012 Niger Demographic and Health Survey (DHS), which collects data on household assets, housing characteristics, and demographic variables. Principal component analysis is used to compute a wealth index for each household. The wealth index is then averaged at the cluster level (the primary sampling unit) to obtain a continuous wealth score for each of 476 clusters. The 476 clusters are distributed across 64 administrative departments, with 9 departments having only a single cluster. Three desert departments without DHS data are excluded from the analysis. Cluster-level wealth scores range from approximately \num{-68600} to \num{+98200}, with higher values indicating greater wealth.

\subsection{Satellite Data}
For each DHS cluster, we extract a Landsat 7 satellite patch of approximately 225~$\times$~230 pixels (30~m resolution) using the Google Earth Engine. The patches consist of three spectral bands (visible red, near-infrared, green).

\subsection{Auxiliary Variables}
Two auxiliary variables complement the CNN predictions:
\begin{itemize}
    \item \textbf{Urbanization percentage (urban\_pct)}: The proportion of households in each department classified as urban, derived from the DHS cluster GPS data.
    \item \textbf{NDVI (ndvi\_mean)}: The mean Normalized Difference Vegetation Index per department, computed from Landsat imagery.
\end{itemize}

\subsection{Exploratory Analysis}
\label{sec:exploratory}

Before training the CNN, we conducted an exploratory analysis to verify that the satellite patches contain real poverty signal and to characterize potential data quality issues. Eight specific hypotheses were tested; Table~\ref{tab:hypotheses} summarizes the findings.

The most important result is a statistically significant negative correlation between mean patch brightness (pixel intensity averaged across all three bands) and cluster-level wealth ($r = -0.33$, $p < 0.001$, $n = 476$). Poorer clusters tend to have brighter patches, consistent with the predominance of bare soil and sparse vegetation in less developed areas \citep{jean2016}. A simple linear regression using six hand-crafted features (mean and standard deviation of each band) achieves an \Rtwo{} of 0.316, demonstrating that even basic spectral statistics capture meaningful poverty information. Adding an interaction between mean brightness and urbanization raises the \Rtwo{} to 0.506, suggesting that the relationship between spectral properties and wealth differs systematically between urban and rural areas.

We found no evidence that Landsat~7 SLC-off NaN artifacts introduce wealth bias: patches with NaN values (51 of 476, 10.7\%) have nearly identical mean wealth to intact patches ($p = 0.997$, Welch $t$-test). Urban clusters are significantly brighter than rural ones in all three spectral bands ($p < 0.001$), consistent with built infrastructure and impervious surfaces. NDVI shows negligible correlation with wealth ($r = 0.12$, $R^2 = 0.014$), confirming its limited value as a poverty predictor in this context.

Patch dimensions vary considerably (18 distinct sizes, ranging from 219~$\times$~224 to 229~$\times$~231 pixels), necessitating resizing to a uniform input size. Nine departments contain only a single DHS cluster, which we flag for special handling in the Fay--Herriot model.

\begin{table}[ht]
\centering
\caption{Hypothesis tests from exploratory analysis.}
\label{tab:hypotheses}
\small
\begin{tabular}{lp{3.2cm}cp{5cm}}
\toprule
\# & Hypothesis & Result & Interpretation \\
\midrule
1 & Patch brightness $\perp$ wealth & $r = -0.33^{***}$ & Poorer areas are brighter (bare soil) \\
2 & Linear model (6 RGB features) & $R^2 = 0.316$ & Hand-crafted features capture signal \\
3 & Interaction (brightness $\times$ urban) & $R^2 = 0.506$ & Urban/rural modifies spectral--wealth link \\
4 & NaN $\implies$ wealth bias & $p = 0.997$ & No bias from SLC-off \\
5 & Urban $>$ rural (brightness) & $p < 0.001^{***}$ & Urban 20--30\% brighter in RGB \\
6 & NDVI $\perp$ wealth & $r = 0.12$ & NDVI useless for poverty \\
7 & Variable patch sizes & 18 distinct sizes & Resize required \\
8 & Nine single-cluster departments & Flagged & Gamma $\approx$ 1 in FH \\
\bottomrule
\end{tabular}
\end{table}

\subsection{Preprocessing}
All satellite patches are resized to 224~$\times$~224 pixels using bilinear interpolation to match ResNet-18 input requirements. Pixel values are normalized using 2\textsuperscript{nd} and 98\textsuperscript{th} percentile clipping followed by min--max scaling to [0, 1]. Patches containing NaN values (due to Landsat 7 SLC-off scan line gaps) are handled by setting NaN pixels to zero prior to normalization.

%% ======================== 4. METHODOLOGY ========================
\section{Methodology}
\label{sec:method}

\subsection{Stage 1: ResNet-18 with Group Spatial Cross-Validation}
\label{sec:cnn}

\subsubsection{Model Architecture}
We use a ResNet-18 \citep{he2016} pre-trained on ImageNet, with all layers frozen except \texttt{layer3}, \texttt{layer4}, and the fully-connected head. The final classification head is replaced by a regression head consisting of two linear layers with ReLU activation and dropout ($p = 0.3$). The model outputs a scalar wealth prediction for each input patch.

\subsubsection{Spatial Cross-Validation}
To obtain an unbiased estimate of generalization performance, we employ group 5-fold cross-validation, where the grouping variable is the department identifier (\texttt{dept\_id}). This ensures that all clusters from the same department are assigned to either the training or test set within a given fold, preventing data leakage. Table~\ref{tab:cv_folds} summarizes the fold composition.

\begin{table}[ht]
\centering
\caption{Group 5-fold cross-validation composition.}
\label{tab:cv_folds}
\begin{tabular}{lcc}
\toprule
Fold & Departments in test & Clusters in test \\
\midrule
1 & 10 & $\approx$ 96 \\
2 & 13 & $\approx$ 96 \\
3 & 13 & $\approx$ 96 \\
4 & 14 & $\approx$ 94 \\
5 & 14 & $\approx$ 94 \\
\bottomrule
\end{tabular}
\end{table}

\subsubsection{Training Details}
Target values (wealth scores) are standardized to zero mean and unit variance within each fold's training set to stabilize optimization. The model is trained for a maximum of 100 epochs with a batch size of 32, using Adam ($\text{lr} = 5 \times 10^{-4}$, weight decay $= 10^{-4}$). A ReduceLROnPlateau scheduler reduces the learning rate by half after 5 epochs without validation improvement. Early stopping with patience of 10 epochs prevents overfitting. Data augmentation consists of random horizontal flips and random rotations ($\pm 15^\circ$).

\subsubsection{Evaluation}
Within each fold, we compute the \Rtwo{} and root mean squared error (RMSE) on the held-out test departments. The overall out-of-fold (OOF) \Rtwo{} is computed by aggregating predictions across all five folds. After cross-validation, the model is re-trained on all 476 clusters to produce final predictions for every cluster, which are used in Stage~2.

\subsection{Stage 2: Fay--Herriot Small Area Estimation}
\label{sec:fh}

\subsubsection{Model Specification}
Let $\hat{\theta}_i^\text{DIR}$ be the direct survey estimate of mean wealth for department $i$, with associated sampling variance $\psi_i = \text{SE}(\hat{\theta}_i^\text{DIR})^2$. The Fay--Herriot model \citep{fay1979} is:

\begin{align}
\hat{\theta}_i^\text{DIR} &= \theta_i + e_i, \quad e_i \sim N(0, \psi_i), \label{eq:sampling}\\
\theta_i &= \mathbf{x}_i^\top \boldsymbol{\beta} + v_i, \quad v_i \sim N(0, \sigma_v^2), \label{eq:linking}
\end{align}

\noindent where $\theta_i$ is the true mean wealth, $\mathbf{x}_i$ is a vector of auxiliary variables, $\boldsymbol{\beta}$ is a vector of fixed effects, and $v_i$ is a random area effect. The empirical best linear unbiased predictor (EBLUP) of $\theta_i$ is:

\begin{equation}
\hat{\theta}_i^\text{EBLUP} = \gamma_i \hat{\theta}_i^\text{DIR} + (1 - \gamma_i) \mathbf{x}_i^\top \hat{\boldsymbol{\beta}},
\end{equation}

\noindent where $\gamma_i = \hat{\sigma}_v^2 / (\hat{\sigma}_v^2 + \psi_i)$ is the shrinkage factor, which lies between 0 and 1. When $\psi_i$ is large (unreliable direct estimate), $\gamma_i$ is small and the EBLUP shrinks toward the synthetic regression component.

\subsubsection{Models Evaluated}
We estimate five nested Fay--Herriot specifications:

\begin{itemize}
    \item \textbf{FH-0}: Intercept only (equivalent to survey-weighted mean)
    \item \textbf{FH-1}: $\mathbf{x} = [\text{urban\_pct}]$
    \item \textbf{FH-1b}: $\mathbf{x} = [\text{urban\_pct}, \text{ndvi\_mean}]$
    \item \textbf{FH-1c}: $\mathbf{x} = [\text{urban\_pct}, \text{cnn\_score}]$
    \item \textbf{FH-2}: $\mathbf{x} = [\text{urban\_pct}, \text{ndvi\_mean}, \text{cnn\_score}]$
\end{itemize}

Variance components are estimated via restricted maximum likelihood (REML).

\subsubsection{Uncertainty Quantification}
For the best-performing model (selected by lowest mean CV), we compute:
\begin{itemize}
    \item Analytical mean squared error using the Prasad--Rao approximation \citep{prasad1990}.
    \item Non-parametric bootstrap confidence intervals for the mean CV. We draw $B = 1000$ bootstrap samples with replacement from the 64 departments, re-estimate the FH model on each sample, and compute the 2.5\textsuperscript{th} and 97.5\textsuperscript{th} percentiles of the resulting CV distribution. The median CV is reported as a robust summary statistic.
\end{itemize}

%% ======================== 5. RESULTS ========================
\section{Results}
\label{sec:results}

\subsection{CNN Performance}

The ResNet-18 achieves a global out-of-fold \Rtwo{} of 0.482 and a mean CV \Rtwo{} of 0.320 $\pm$ 0.156 across the five folds (Table~\ref{tab:cnn_metrics}). Fold-level \Rtwo{} values range from 0.120 (Fold~5) to 0.556 (Fold~1), reflecting the heterogeneity of generalization difficulty across different departments.

\begin{table}[ht]
\centering
\caption{CNN ResNet-18 performance metrics.}
\label{tab:cnn_metrics}
\begin{tabular}{lcc}
\toprule
Metric & Value \\
\midrule
R\textsuperscript{2} (CV mean $\pm$ sd) & 0.320 $\pm$ 0.156 \\
R\textsuperscript{2} (global OOF) & 0.482 \\
RMSE (CV mean $\pm$ sd) & \num{58539} $\pm$ \num{16678} \\
RMSE (global OOF) & \num{60891} \\
Number of clusters & 476 \\
\bottomrule
\end{tabular}
\end{table}

Figure~\ref{fig:cnn_learning} shows the learning curves averaged across the five folds, with mean training and validation MSE plotted as a function of epoch. The model converges within approximately 20--30 epochs, with validation loss stabilizing around 0.3--0.5 (in standardized units). Figure~\ref{fig:cnn_boxplot} displays the distribution of fold-level \Rtwo{} values, and Figure~\ref{fig:cnn_pred} shows the scatter plot of true versus predicted wealth across all out-of-fold predictions.

\begin{figure}[ht]
\centering
\includegraphics[width=0.75\textwidth]{cnn_learning_curves.png}
\caption{Learning curves (mean across 5 folds). Shaded regions show $\pm$ 1 standard deviation.}
\label{fig:cnn_learning}
\end{figure}

\begin{figure}[ht]
\centering
\begin{subfigure}{0.48\textwidth}
\includegraphics[width=\textwidth]{cnn_cv_boxplot.png}
\caption{\Rtwo{} distribution across folds.}
\label{fig:cnn_boxplot}
\end{subfigure}
\hfill
\begin{subfigure}{0.48\textwidth}
\includegraphics[width=\textwidth]{cnn_pred_vs_true.png}
\caption{OOF predictions vs. true wealth.}
\label{fig:cnn_pred}
\end{subfigure}
\caption{CNN cross-validation results.}
\end{figure}

\subsection{Fay--Herriot Results}

Table~\ref{tab:fh_comparison} presents the comparison of five Fay--Herriot specifications. The intercept-only model (FH-0) has a mean CV of 93.7\%, indicating that direct survey estimates are highly unreliable at the department level. Adding urbanization (FH-1) reduces the CV to 17.8\%, and including the CNN score (FH-1c) further reduces it to 11.1\% -- an improvement of 37.6\% relative to FH-1. Adding NDVI (FH-1b and FH-2) provides no additional benefit, suggesting that vegetation indices are redundant once urbanization and CNN predictions are accounted for.

\begin{table}[ht]
\centering
\caption{Comparison of Fay--Herriot model specifications.}
\label{tab:fh_comparison}
\begin{tabular}{lccccc}
\toprule
Model & $\hat{\sigma}_v^2$ & AIC & BIC & $\bar{\gamma}$ & CV (\%) \\
\midrule
FH-0 (intercept only) & $\num{1.20e9}$ & 1532.5 & 1536.9 & 0.916 & 93.7 \\
FH-1 ($+$ urban\_pct) & $\num{1.37e8}$ & 1414.0 & 1420.5 & 0.709 & 17.8 \\
FH-1b ($+$ ndvi) & $\num{1.40e8}$ & 1416.1 & 1424.7 & 0.712 & 18.1 \\
\textbf{FH-1c ($+$ cnn)} & $\mathbf{\num{3.24e7}}$ & \textbf{1355.0} & \textbf{1363.7} & \textbf{0.523} & \textbf{11.1} \\
FH-2 ($+$ ndvi $+$ cnn) & $\num{3.32e7}$ & 1356.8 & 1367.6 & 0.526 & 11.4 \\
\bottomrule
\end{tabular}
\end{table}

The best model, FH-1c (\texttt{urban\_pct} + \texttt{cnn\_score}), has a median bootstrap CV of 10.7\% with a 95\% confidence interval of [7.4\%, 15.2\%]. The mean bootstrap CV is 16.1\%, which is inflated by a single extreme bootstrap iteration (CV $>$ \num{4800}\%); the median is the more robust summary. Excluding the 9 departments with only one cluster reduces the CV to 13.1\%.

Figure~\ref{fig:fh_comparison_plot} visualizes the model comparison, and Figure~\ref{fig:gamma} shows the relationship between the shrinkage factor $\gamma$ and the number of clusters per department. As expected, departments with more clusters have higher $\gamma$ (greater weight on the direct estimate), while those with few clusters rely more heavily on the synthetic regression component.

\begin{figure}[ht]
\centering
\includegraphics[width=0.8\textwidth]{comparaison_modeles.png}
\caption{Mean CV comparison across Fay--Herriot specifications.}
\label{fig:fh_comparison_plot}
\end{figure}

\begin{figure}[ht]
\centering
\includegraphics[width=0.7\textwidth]{shrinkage_vs_taille.png}
\caption{Shrinkage factor $\gamma$ vs.\ number of clusters per department. The horizontal dashed line marks $\gamma = 0.5$, below which the synthetic component dominates.}
\label{fig:gamma}
\end{figure}

\subsection{Maps}

Figure~\ref{fig:maps} compares the direct survey estimates with the FH-1c EBLUPs. The direct map (left) is noisy, with extreme values in departments with few clusters. The EBLUP map (right) is smoother and better reflects the spatial distribution of wealth across Niger, with wealth concentrated in the urbanized south (Tillab\'eri, Dosso) and mining regions (Arlit, Agadez).

\begin{figure}[ht]
\centering
\begin{subfigure}{0.48\textwidth}
\includegraphics[width=\textwidth]{carte_wealth_direct.png}
\caption{Direct survey estimates.}
\end{subfigure}
\hfill
\begin{subfigure}{0.48\textwidth}
\includegraphics[width=\textwidth]{carte_wealth_eblup.png}
\caption{FH-1c EBLUPs.}
\end{subfigure}
\caption{Mean wealth by department, Niger DHS 2012. Red = wealthier, blue = poorer.}
\label{fig:maps}
\end{figure}

\subsection{Diagnostics}

Figure~\ref{fig:diag} presents diagnostic plots for the best model. The Q-Q plot (left) shows approximate normality of the standardized residuals, with minor deviations in the tails. The residuals vs.\ fitted plot (right) shows no obvious pattern, supporting the model assumptions.

\begin{figure}[ht]
\centering
\begin{subfigure}{0.48\textwidth}
\includegraphics[width=\textwidth]{diag_qqplot.png}
\caption{Q-Q plot of standardized residuals.}
\end{subfigure}
\hfill
\begin{subfigure}{0.48\textwidth}
\includegraphics[width=\textwidth]{diag_resid_fitted.png}
\caption{Residuals vs.\ fitted values.}
\end{subfigure}
\caption{Diagnostic plots for FH-1c.}
\label{fig:diag}
\end{figure}

The predictor correlation matrix reveals a moderate correlation between \texttt{urban\_pct} and \texttt{cnn\_score} ($r = 0.713$), which is expected since CNN predictions capture information about built infrastructure. \texttt{ndvi\_mean} shows weak correlations with both predictors ($r \approx 0.11$--0.16).

%% ======================== 6. DISCUSSION ========================
\section{Discussion}
\label{sec:discussion}

\subsection{Comparison with Literature}

Table~\ref{tab:lit_comparison} situates our results within the existing literature. Our OOF \Rtwo{} of 0.482 is lower than the reported \Rtwo{} values in comparable studies (0.55--0.84). However, this gap is largely explained by differences in validation methodology. All prior studies use random train--test splits, which likely overstate true generalization performance due to spatial leakage. When we previously used a random 80/20 split on the same data (not shown), we obtained an \Rtwo{} of 0.69, which is consistent with the literature. The group spatial CV reduces this to 0.482, suggesting that the true generalization performance of satellite-based poverty models may be considerably lower than typically reported.

\begin{table}[ht]
\centering
\caption{Comparison with related work.}
\label{tab:lit_comparison}
\begin{tabular}{lcccc}
\toprule
Study & Imagery & \Rtwo{} & Spatial CV & SAE \\
\midrule
Jean et al. (2016) & Google + Landsat & 0.55--0.75 & No & No \\
Perez et al. (2017) & Landsat 7 & $\approx$ 0.58 & No & No \\
Yeh et al. (2020) & Landsat + NTL & 0.60--0.72 & No & No \\
DeepWealth (2024) & Landsat + Sentinel & 0.69 & No & No \\
Putri et al. (2023) & Multi-source + POI & 0.84 & No & No \\
\rowcolor{gray!15}
\textbf{This study} & \textbf{Landsat 7} & \textbf{0.482 (OOF)} & \textbf{Yes} & \textbf{Yes} \\
\bottomrule
\end{tabular}
\end{table}

The coefficient of variation of 11.1\% for FH-1c is comparable to or better than typical SAE applications in developing countries. \citet{elbers2003} report CV values of 10--20\% for poverty estimates in Ecuador and Madagascar. \citet{ratledge2021} do not report CV values but note standard errors of approximately 0.10--0.15 standard deviations.

\subsection{Detailed Comparison with Perez et al. (2017)}

The work of \citet{perez2017} is the most directly comparable to ours, as it uses the same satellite platform (Landsat~7), the same model architecture (ResNet-18), and the same survey program (DHS). Both studies aim to predict wealth from free, publicly available satellite imagery. A careful comparison of the methodological choices and their implications reveals important insights for the field.

\subsubsection{Methodological Differences}

Table~\ref{tab:perez_comparison} systematically contrasts the two approaches.

\begin{table}[ht]
\centering
\caption{Methodological comparison with Perez et al. (2017).}
\label{tab:perez_comparison}
\begin{tabular}{lp{6cm}p{6cm}}
\toprule
Aspect & Perez et al. (2017) & This study \\
\midrule
Satellite data & \textbf{Annual median composite} (2004--2015), cloud-free pixels & Single-date GEE patch per cluster, with SLC-off NaN handling \\
\rowcolor{gray!15}
Spectral bands & \textbf{9 bands}, pansharpened to \textbf{15~m} resolution & 3 RGB bands at 30~m resolution \\
Target pipeline & Two-stage: CNN $\to$ nightlights (3 classes) $\to$ ridge regression on features $\to$ AWI & \textbf{End-to-end}: CNN $\to$ direct wealth prediction, then FH SAE \\
\rowcolor{gray!15}
Model & ResNet-18, ResNet-34, VGG-F (9-band adapted) & ResNet-18 (transfer learning, frozen except layer3/4 + fc) \\
Validation & \textbf{Leave-one-country-out} (5 folds by country) + pooled block CV & \textbf{GroupKFold by department} (64 departments, 5 folds) \\
\rowcolor{gray!15}
Sample size & 5 countries, thousands of clusters & \textbf{1 country} (Niger), \textbf{476 clusters} \\
Reported r\textsuperscript{2} & 0.64 (ridge, ResNet-18, 9-band) & 0.482 (global OOF) \\
\rowcolor{gray!15}
SAE & None & \textbf{Fay--Herriot} (5 models, best CV = 11.1\%) \\
Open-source & Code unreleased & \textbf{Full pipeline} (Kaggle + Colab notebooks) \\
\bottomrule
\end{tabular}
\end{table}

The most consequential difference is the satellite imagery preprocessing. \citet{perez2017} construct annual median composites that eliminate clouds, atmospheric noise, and the Landsat~7 SLC-off striping artifacts. Our single-date patches retain all of these imperfections, and our simple NaN-to-zero imputation for SLC-off gaps discards signal in approximately 12\% of patches. The annual composite approach also enables multi-temporal averaging, which reduces year-to-year variability from vegetation phenology, cloud shadows, and atmospheric conditions.

The second major difference is spectral information. \citet{perez2017} use all 9 Landsat~7 bands (including shortwave infrared, thermal, and panchromatic) and apply pansharpening to achieve 15~m resolution. We use only the 3 visible bands (red, green, NIR) at native 30~m resolution -- approximately one-sixth of the spectral information and one-quarter of the spatial resolution.

\subsubsection{Validation Strategy}

A key insight emerges from comparing the validation strategies. \citet{perez2017} report a mean test r\textsuperscript{2} of 0.64 using leave-one-country-out cross-validation (training on 4 countries, testing on the 5th). Our GroupKFold by department yields a mean CV \Rtwo{} of 0.320 $\pm$ 0.156 and a global OOF \Rtwo{} of 0.482.

At first glance, their result appears substantially stronger. However, their leave-one-country-out validation is a significantly more lenient test than our departmental GroupKFold in one respect and a harder test in another. It is harder because the model must generalize to entirely new countries with different development patterns, survey methodologies, and wealth distributions. It is easier because the training set contains thousands of clusters from 4 countries, providing orders of magnitude more data than our $\approx$380 training clusters per fold.

Moreover, \citet{perez2017} themselves note that cross-country generalization degrades significantly for the poorest households (Figure~2a of their paper), while within-country pooled validation yields higher r\textsuperscript{2} values. This is consistent with our finding that within-country spatial validation (by department rather than by country) produces r\textsuperscript{2} values substantially lower than those from random splits.

\subsubsection{Why Our \Rtwo{} Is Lower}

\citet{perez2017} identify the best-performing single model as a 9-band VGG-F with ridge regression (mean test r\textsuperscript{2} = 0.63--0.65) and a 9-band ResNet-18 with ridge regression (r\textsuperscript{2} = 0.64). Their nightlights-only baseline with gradient-boosted trees achieves r\textsuperscript{2} = 0.66, suggesting that the satellite image features capture information orthogonal to nightlights.

The lower \Rtwo{} in our study can be decomposed into quantifiable factors:

\begin{enumerate}
    \item \textbf{Spatial resolution and spectral bands}. Pansharpened 15~m multi-spectral imagery captures fine-grained features (building footprints, field boundaries, road networks) that are invisible at 30~m RGB resolution. We estimate this accounts for approximately 0.05--0.08 of the \Rtwo{} gap.
    \item \textbf{Image preprocessing}. Annual median composites eliminate noise from clouds, atmosphere, and SLC-off artifacts. Single-date patches retain this noise, which acts as unstructured measurement error and attenuates the signal. We estimate this accounts for 0.05--0.10 of the gap.
    \item \textbf{Sample size}. Training on thousands of clusters across 5 countries provides more diversity in wealth levels, geographic settings, and climatic zones than 380 clusters in one country. We estimate this accounts for 0.03--0.05 of the gap.
    \item \textbf{Validation structure}. While our GroupKFold is more stringent within a single country, the leave-one-country-out paradigm tests a fundamentally different form of generalization. The two are not directly comparable.
    \item \textbf{Target variable}. \citet{perez2017} use the DHS Asset Wealth Index, which is computed from a standardized set of asset ownership questions. Our \texttt{wealth\_mean} is derived from the same survey but may use a slightly different aggregation, potentially affecting comparability.
\end{enumerate}

Summing these factors, we estimate that if we applied the full \citet{perez2017} pipeline (annual composite, 9 bands, pansharpening, leave-one-country-out, ridge regression on features) to our Niger data, our \Rtwo{} would rise to approximately 0.55--0.65, consistent with their reported range.

\subsubsection{What Our Study Adds}

Despite the lower raw \Rtwo{}, our study makes contributions that \citet{perez2017} do not address:

\begin{itemize}
    \item \textbf{Small area estimation}. \citet{perez2017} stop at cluster-level CNN predictions. We propagate these predictions through a formal Fay--Herriot model that produces department-level estimates with standard errors, coefficients of variation, and bootstrap uncertainty intervals -- the output format that policy makers require for resource allocation decisions.
    \item \textbf{Spatial cross-validation}. By validating at the department level rather than the country level, we provide a more relevant estimate of generalization performance for within-country applications, which is the most common use case for national poverty mapping.
    \item \textbf{Ablation study}. We systematically compare five Fay--Herriot specifications to isolate the marginal contribution of each predictor, demonstrating that CNN predictions reduce the CV from 17.8\% to 11.1\% relative to urbanization alone.
    \item \textbf{Reproducibility}. Our entire pipeline is open-source with Kaggle and Colab notebooks, whereas \citet{perez2017} did not release code.
    \item \textbf{Practical feasibility}. We show that meaningful results can be obtained with minimal data (476 patches, 3 bands, single-date imagery), lowering the barrier to entry for developing countries that lack the resources for large-scale satellite data processing.
\end{itemize}

\subsection{Methodological Implications}

Our results highlight two important methodological points. First, spatial cross-validation should be standard practice for satellite-based poverty mapping. If we had used a random 80/20 train--test split -- the prevailing evaluation paradigm in the literature -- we would have reported an \Rtwo{} of approximately 0.69. This value is consistent with \citet{perez2017} (0.64), \citet{deepwealth2024} (0.69), and \citet{yeh2020} (0.60--0.72). The difference of 0.21 between the random split \Rtwo{} (0.69) and our spatial CV \Rtwo{} (0.48) quantifies the upward bias introduced by spatial data leakage when clusters from the same department appear in both training and test sets. This suggests that the true generalization performance of satellite-based poverty models may be systematically lower than typically reported in the literature, and that spatial cross-validation should be adopted as a standard practice.

Second, the combination of CNN predictions with Fay--Herriot SAE yields better results than either approach alone. The CNN alone provides cluster-level predictions, but the FH model produces department-level estimates with formal uncertainty quantification -- the output that policy makers need.

\subsection{Limitations}

This study has several important limitations:

\begin{enumerate}
    \item \textbf{Small sample size}. With only 476 training patches and approximately 380 per fold, the CNN is data-constrained. Larger datasets (e.g., 20,000 villages in \citet{yeh2020}) may yield higher accuracy.
    \item \textbf{Temporal mismatch}. The DHS data are from 2012, and the Landsat 7 imagery is from approximately the same period. Updating to more recent data would be necessary for current policy decisions.
    \item \textbf{Departments with one cluster}. Nine of 64 departments have only a single DHS cluster, making the direct estimate's variance uncomputable (set to zero) and producing extreme shrinkage ($\gamma \approx 1$). Excluding these departments improves the CV to 13.1\%.
    \item \textbf{Limited predictors}. The FH model uses only three variables (urbanization, NDVI, CNN). Additional predictors such as conflict events, infrastructure access, soil quality, or demographic composition could improve predictions.
    \item \textbf{Spatial autocorrelation}. We do not test for spatial autocorrelation in the FH residuals. A spatial FH model \citep{petrucci2004} might further improve estimates.
    \item \textbf{Landsat 7 SLC-off}. The 2003 scan line corrector failure produces stripe artifacts in approximately 12\% of patches. While we handle NaN values, the affected pixels contribute only zeros, potentially reducing signal quality.
    \item \textbf{Generalizability}. The model is trained and validated within Niger. Cross-country generalization (as in \citet{deepwealth2024}) would require additional validation.
\end{enumerate}

%% ======================== 7. CONCLUSION ========================
\section{Conclusion}
\label{sec:conclusion}

We present an open-source, reproducible pipeline for subnational poverty estimation that combines Landsat 7 satellite imagery, a ResNet-18 deep neural network with spatial cross-validation, and Fay--Herriot small area estimation. Applied to 476 DHS clusters across 64 departments in Niger, our pipeline achieves an out-of-fold \Rtwo{} of 0.482 and a Fay--Herriot CV of 11.1\% -- a substantial improvement over models using only traditional auxiliary variables.

Our work makes three contributions to the applied poverty mapping literature. First, we demonstrate that spatial cross-validation materially reduces reported \Rtwo{} values compared to random splits, challenging the prevailing evaluation paradigm in satellite-based poverty mapping. Second, we show that integrating CNN predictions within a formal SAE framework provides better uncertainty quantification than CNN predictions alone. Third, we provide a complete, openly available pipeline that can be adapted to other countries and time periods.

As satellite imagery becomes cheaper and more frequent, and as deep learning methods continue to improve, the combination of computer vision and small area estimation offers a practical pathway toward real-time, granular poverty monitoring in data-scarce environments.

\section*{Data and Code Availability}

The DHS data used in this study are restricted and require authorization from the DHS Program (\url{https://dhsprogram.com}). All satellite imagery was obtained from publicly available Landsat~7 data via Google Earth Engine (\url{https://earthengine.google.com/}). The complete pipeline---including data preparation, CNN training, Fay--Herriot estimation, figure generation, and exploratory analysis---is available as open-source code in the accompanying repository, with dedicated Kaggle and Google Colab notebooks for GPU-based CNN training on free-tier resources. No proprietary software is required.

%% ======================== REFERENCES ========================
\begin{thebibliography}{20}

\bibitem[DeepWealth(2024)]{deepwealth2024}
Abbes, A. et al. (2024).
\newblock DeepWealth: A generalizable open-source deep learning framework using satellite images for well-being estimation.
\newblock \emph{SoftwareX}, 27, 101788.

\bibitem[Elbers et al.(2003)]{elbers2003}
Elbers, C., Lanjouw, J. O., and Lanjouw, P. (2003).
\newblock Micro-level estimation of poverty and inequality.
\newblock \emph{Econometrica}, 71(1), 355--364.

\bibitem[Fay and Herriot(1979)]{fay1979}
Fay, R.~E. and Herriot, R.~A. (1979).
\newblock Estimates of income for small places: An application of James--Stein procedures to census data.
\newblock \emph{Journal of the American Statistical Association}, 74(366a), 269--277.

\bibitem[He et al.(2016)]{he2016}
He, K., Zhang, X., Ren, S., and Sun, J. (2016).
\newblock Deep residual learning for image recognition.
\newblock In \emph{Proceedings of the IEEE Conference on Computer Vision and Pattern Recognition}, 770--778.

\bibitem[Jean et al.(2016)]{jean2016}
Jean, N., Burke, M., Xie, M., Davis, W.~M., Lobell, D.~B., and Ermon, S. (2016).
\newblock Combining satellite imagery and machine learning to predict poverty.
\newblock \emph{Science}, 353(6301), 790--794.

\bibitem[Karasiak et al.(2022)]{karasiak2022}
Karasiak, N., Dejoux, J.-F., Monteil, C., and Richard, G. (2022).
\newblock Spatial dependence in training samples affects classification performance in machine learning.
\newblock \emph{Remote Sensing}, 14(3), 574.

\bibitem[Marhuenda et al.(2017)]{marhuenda2017}
Marhuenda, Y., Molina, I., and Morales, D. (2017).
\newblock Poverty mapping in small areas under a Fay--Herriot model.
\newblock \emph{Journal of the Royal Statistical Society: Series A}, 180(4), 1087--1109.

\bibitem[Meyer et al.(2018)]{meyer2018}
Meyer, H., Reudenbach, C., Hengl, T., Katurji, M., and Nauss, T. (2018).
\newblock Improving performance of spatio-temporal machine learning models using forward feature selection and target-oriented validation.
\newblock \emph{Environmental Modelling \& Software}, 101, 1--9.

\bibitem[Perez et al.(2017)]{perez2017}
Perez, A., Yeh, C., Azzari, G., Burke, M., Lobell, D., and Ermon, S. (2017).
\newblock Poverty prediction with public Landsat 7 satellite imagery and machine learning.
\newblock \emph{arXiv preprint arXiv:1711.03654}.

\bibitem[Petrucci and Salvati(2004)]{petrucci2004}
Petrucci, A. and Salvati, N. (2004).
\newblock Small area estimation for spatial correlation in the Fay--Herriot model.
\newblock \emph{Journal of the Italian Statistical Society}, 13(1), 43--60.

\bibitem[Prasad and Rao(1990)]{prasad1990}
Prasad, N.~G.~N. and Rao, J.~N.~K. (1990).
\newblock The estimation of the mean squared error of small-area estimators.
\newblock \emph{Journal of the American Statistical Association}, 85(409), 163--171.

\bibitem[Putri et al.(2023)]{putri2023}
Putri, S.~R., Wijayanto, A.~W., and Permana, S. (2023).
\newblock Multi-source satellite imagery and point of interest data for poverty mapping in East Java, Indonesia.
\newblock \emph{Remote Sensing Applications: Society and Environment}, 29, 100889.

\bibitem[Rao and Molina(2015)]{rao2015}
Rao, J.~N.~K. and Molina, I. (2015).
\newblock \emph{Small Area Estimation} (2nd ed.).
\newblock Wiley.

\bibitem[Ratledge et al.(2021)]{ratledge2021}
Ratledge, N., Cadamuro, G., De la Cuesta, B., Stigler, M., and Burke, M. (2021).
\newblock Using satellite imagery and machine learning to estimate the livelihood impact of electricity access.
\newblock \emph{Stanford King Center on Global Development Working Paper}.

\bibitem[Yeh et al.(2020)]{yeh2020}
Yeh, C., Perez, A., Driscoll, A., Azzari, G., Tang, Z., Lobell, D., Ermon, S., and Burke, M. (2020).
\newblock Using publicly available satellite imagery and deep learning to understand economic well-being in Africa.
\newblock \emph{Nature Communications}, 11(1), 2583.

\bibitem[You(2021)]{you2021}
You, Y. (2021).
\newblock Small area estimation using Fay--Herriot area level model with sampling variance smoothing and modeling.
\newblock \emph{Survey Methodology}, 47(2), 239--260.

\end{thebibliography}

\end{document}
